% =====================================================================
% 2. Evaluation framework (PRDL, Maria Martinez)
% Own block in depth: evaluation protocol (T7; DECISIONS D3, D4, D14;
% block note 03). Shared training loop (T8, OF + MM).
% =====================================================================
\section{Evaluation framework}\label{sec:evaluation}

I designed and wrote the evaluation module (task T7, reviewed by Oriana); the training loop was written
by the two of us (T8). The aim of the protocol is that every number in this report means the same thing
for every model, and that a difference is claimed only when it is larger than what chance alone would
produce.

\subsection{Loss and outputs}

Each model returns five logits, one per superclass, with no sigmoid inside the model. The loss is the
binary cross-entropy with logits, averaged over the five classes and the batch: each class is an
independent yes/no question, which is what a multi-label task needs (a softmax would force the five
probabilities to compete and sum to 1). Probabilities are obtained afterwards as
$\sigma(\text{logit})$. The optimiser is Adam with learning rate $10^{-3}$, whose hyper-parameters
``typically require little tuning'' \cite[p.~1]{kingma2015adam}; batch 64, at most 50 epochs, and early
stopping with a patience of 5 epochs on the validation macro AUROC, after which the weights of the best
epoch are restored.

\subsection{Metrics}

\paragraph{Macro AUROC.} The main metric is the area under the ROC curve of each class, averaged over the
five classes. The AUROC is the probability that a randomly chosen positive record gets a higher score than
a randomly chosen negative one, so it needs no threshold and does not change with the class prevalence.
Averaging over classes (macro) gives HYP, the rarest class, the same weight as NORM. It is also the
metric of the benchmark, ``obtained by averaging class-wise AUCs over all classes''
\cite[p.~1522]{strodthoff2021deep}, which makes a comparison possible (Section~\ref{sec:discussion}).

\paragraph{Macro F1.} As a secondary metric, which does depend on a decision, each class gets the
threshold that maximises its F1 on the validation fold (grid 0.01 to 0.99); the same thresholds are
applied unchanged to the test fold. A fixed threshold of 0.5 for every class was rejected because it
penalises the rare classes, whose probabilities rarely reach 0.5.

\subsection{Uncertainty and comparison of models}

\paragraph{Patient bootstrap.} The test fold is a sample, so every metric has a confidence interval: the
predictions are resampled 1000 times and the 2.5 and 97.5 percentiles of the metric are reported. The
resampling unit is the patient, not the record: when a patient is drawn, all of his or her records enter
together. Records of the same patient are correlated, and resampling them one by one would pretend that
we have more independent evidence than we do, giving intervals that are too narrow. The model is not
retrained, only its predictions are resampled.

\paragraph{Paired comparison.} Two models evaluated on the same records can be compared on the same
resamples: the draws depend only on the patients and on the seed, so for each resample we compute the
difference of their macro AUROCs and take its 95\,\% interval. If the interval contains 0, neither model
is declared better. Pairing removes the variation that comes from which patients happen to be in the
sample, so it detects much smaller differences than comparing two separate intervals. Every model saves
its predictions record by record, so any two runs can be compared afterwards.

\paragraph{Seeds.} A single training run is itself a random draw (initial weights, batch order,
dropout). Every configuration was trained with three seeds (42, 43, 44), and a configuration is compared
with its reference trained with the same seed. Tables report the mean and standard deviation over the
seeds and in how many seeds the paired interval excludes 0 in each direction. A factor is said to change
the result only when it does so consistently, that is, in all three seeds; one or two seeds of three
are reported as weak evidence. The bootstrap interval reflects the finite test set; the seeds reflect the
randomness of training.

\subsection{Selection and use of the test fold}

All ablations and every choice of a final model use the validation fold only. The selection rule of each
block (the configuration with the highest mean validation macro AUROC; for RNN + attention, the best of
the attention configurations) was written in the decisions log before the runs of that block, and each
final model was trained again with the three seeds and evaluated once on the test fold; the training code
refuses a second test run of the same model and seed. A rule fixed in advance can select a configuration
that is only marginally better than its baseline; we keep its choice and say so, because changing it
after seeing the test results would be selection on noise.

\subsection{Traceability}

Every run writes a log with its configuration, seed, data description, learning curves, metrics,
software versions and code commit, and its predictions. No number of this report is typed by hand: each
is a macro generated from those logs by a script, and only the logs listed in a fixed file can appear in a
report, so an old or test run cannot slip in. The only numbers typed by hand in the whole project are the
published benchmark values, each stored with its file, page and quote.

As a floor for every comparison, a trivial model (logistic regression on the mean and standard deviation
of each lead, 24 numbers per record) reaches a test macro AUROC of \trivialTestMacroAuroc{}
[\trivialTestMacroAurocLo{}, \trivialTestMacroAurocHi{}]: any real architecture must beat it clearly.
